\documentclass{situationist-preprint}

\title{Understanding Without Knowledge}
\msid{EPI-24-a}
\msfamily{EPI-24}
\msdiscipline{Epistemology}
\mstype{Principal manuscript}
\mskeywords{understanding; knowledge; explanation; grasp; epistemic luck; safety; testimony; idealization}

\begin{document}
\maketitle

\begin{abstract}
Is understanding a species of knowledge? On the reductive view, to understand why something is the case is to know why it is the case. We argue that understanding is not a kind of knowledge, because the two come apart in both directions. We characterize understanding why $p$ as the grasp of an explanation of $p$ that is correct in its core: the explanation is true enough in what it says about the dependence of $p$ on its factors, and the subject can use it to draw correct conclusions about how $p$ would have differed had its explanatory factors differed. Knowledge, by contrast, is safe, justified true belief. We prove that a belief can fail to be safe while the grasp it supports is perfectly in order, so that understanding survives the kinds of luck that defeat knowledge; that understanding tolerates falsehoods in the periphery of the explanation and in its idealizations, where knowledge of the corresponding propositions is impossible; and that knowledge of a correct explanation, acquired by testimony, does not bring the grasp that understanding requires. The two states answer to different epistemic ends: knowledge to getting things right non-accidentally, understanding to seeing how things hang together.
\end{abstract}

\section{The problem}

Much of what is valued in inquiry is not the accumulation of truths but understanding: seeing why the tides follow the moon, why the Weimar Republic fell, why a proof works. Philosophers who took knowledge to be the central epistemic notion assumed, mostly without argument, that understanding is a kind of knowledge, namely knowledge of causes or explanations. The assumption has an ancient pedigree, since Aristotle's \emph{epist\=em\=e} is knowledge of a thing through its cause, and it was natural to translate this both as \enquote{scientific knowledge} and as \enquote{understanding} \citep{Burnyeat1981}.

The assumption was challenged at the start of the century. \citet{Zagzebski2001} argued that understanding, unlike knowledge, concerns the grasp of structure and is directed at systems of propositions rather than at single ones. \citet{Kvanvig2003} argued that understanding is not even factive in the way knowledge is, and that it is compatible with the kind of luck that Gettier cases exhibit. His example is a student who learns the history of the Comanche from a book that is, unknown to her, one of very few accurate books in a collection of fakes; she understands the history, he holds, though her beliefs are unsafe. \citet{Elgin2007,Elgin2017} argued that much scientific understanding is achieved through models and idealizations known to be false, and that understanding tolerates \enquote{felicitous falsehoods}. \citet{Pritchard2010} took an intermediate position: understanding is compatible with environmental luck, of the kind exhibited by the fake-barn case, but not with the intervening luck of Gettier's own cases; and it is not entailed by knowledge, since one can know why something happened on testimony without grasping the explanation. \citet{Hills2016} gave an account of understanding why in terms of abilities to follow, give and apply explanations, and argued on that basis that understanding and knowledge come apart.

The reductive view has been defended in response. \citet{Grimm2006} argued that Kvanvig's case shows at most that understanding is compatible with some luck, as knowledge itself is on some accounts, and that understanding why is a species of knowledge why; \citet{Sliwa2015} argued that understanding why $p$ is knowing why $p$, and that the degrees of understanding are degrees of knowledge of the relevant explanatory facts; \citet{Khalifa2017} argued that understanding is a matter of having scientific knowledge of explanations, so that whatever distinguishes understanding from knowledge belongs to the knowledge of more explanatory propositions.

We side with the non-reductive view, in a form that is stronger than Pritchard's and more precise than Kvanvig's. With the analysis of knowledge as safe, justified true belief (EPI-01), we can say exactly which conditions on knowledge understanding lacks and which conditions on understanding knowledge lacks.

\section{Framework}

\subsection{Knowledge}

We adopt the analysis of EPI-01.

\begin{definition}[Knowledge]\label{d:know}
$S$ knows that $p$ if and only if $S$ believes that $p$ on some basis, is justified in so believing, and the belief is safe: in every close world in which $S$ believes that $p$ on the same basis, $p$ is true.
\end{definition}

Safety implies truth, since the actual world is close to itself. A justified true belief is not knowledge exactly when there is a close world in which the subject believes the same proposition on the same basis and the proposition is false.

\subsection{Explanations and grasp}

We follow SCI-03 in taking an explanation to locate its explanandum in a pattern of dependence: it says on what factors the explanandum depends and how. An explanation of why $p$ can be represented as a set $E$ of propositions, which we divide into a \emph{core}, the propositions stating the factors on which $p$ depends and the form of the dependence, and a \emph{periphery}, the propositions that serve to apply or simplify the core, including idealizing assumptions and approximate values.

\begin{definition}[True enough]\label{d:trueenough}
A representation $E$ of an explanation of $p$ is \emph{true enough} if its core propositions are true and its peripheral propositions, where false, depart from the truth only in ways that do not change which factors $p$ depends on or the direction of the dependence.
\end{definition}

The definition makes precise Elgin's notion of a felicitous falsehood. The ideal gas law is false of every actual gas, since actual molecules have volume and attract one another; but in explaining why the pressure of a gas rises when it is heated at constant volume, the idealization departs from the truth in a way that leaves the dependence of pressure on temperature, and its direction, as they are.

\begin{definition}[Grasp]\label{d:grasp}
$S$ \emph{grasps} an explanation $E$ of $p$ if $S$ can use $E$ to answer correctly a range of what-if-things-had-been-different questions about $p$: for the factors on which, according to $E$, $p$ depends, $S$ can say how $p$ would have differed had they differed, and can apply $E$ to relevantly similar cases.
\end{definition}

The notion of what-if-things-had-been-different questions is Woodward's \citep{Woodward2003}; the abilities are close to those Hills lists \citep{Hills2016}. Grasp is an ability, and abilities are not beliefs. A subject may believe every proposition in $E$ and lack the ability, as a student who has memorized a proof without being able to say what would fail if a premise were dropped. Grasp comes in degrees, measured by the range of questions the subject can answer, and so therefore does understanding.

\begin{definition}[Understanding why]\label{d:understand}
$S$ \emph{understands why} $p$ if and only if $p$ is true and $S$ accepts a representation $E$ of an explanation of $p$ that is true enough and that $S$ grasps.
\end{definition}

We use \enquote{accepts} rather than \enquote{believes} for the attitude toward $E$, since a scientist who explains by an idealized model need not believe the idealization; she uses it as a premise in reasoning about the phenomenon \citep{Elgin2017}. Toward the core of $E$ the attitude is belief.

Definition~\ref{d:understand} is deliberately austere. It contains no safety condition and no justification condition. We now argue that this is correct, and that the two conditions it omits are exactly those that separate understanding from knowledge.

\section{Understanding without knowledge}

\begin{principle}[Grasp is basis-independent]\label{p:basis}
Whether $S$ grasps $E$ depends on the abilities $S$ has with respect to $E$, and not on the reliability of the source from which $S$ acquired $E$.
\end{principle}

Principle~\ref{p:basis} follows from the conception of grasp as an ability. The ability to say how the pressure of a gas would change if its volume were halved is exercised in reasoning from the explanation, and it is the same ability whether the explanation was learned from a reliable textbook or from an unreliable one that happens to be correct. What the source affects is whether the subject's acceptance of $E$ could easily have been acceptance of a false explanation; it does not affect what she can do with $E$ once she has it.

\begin{theorem}\label{t:luck}
Understanding why $p$ does not entail knowledge of the core propositions of the explanation understood.
\end{theorem}

\begin{proof}
Consider a student who learns why the Comanche rose to dominance on the southern plains from a book that is accurate, in a library in which most books on the subject are fabrications that she cannot distinguish from it. She believes the core propositions of the book's explanation (the acquisition of horses, the shift to bison hunting, the control of trade routes), they are true, and she has worked through the explanation until she can say what would have happened had the horse arrived later or the trade routes been held by others. By Definition~\ref{d:understand} and Principle~\ref{p:basis}, she understands why the Comanche rose to dominance. But there are close worlds in which she takes down a fabricated book, on the same basis (trust in an apparently scholarly history from that library), and believes a false explanation. Her beliefs in the core propositions are therefore unsafe, and by Definition~\ref{d:know} she does not know them.
\end{proof}

The case is Kvanvig's, with the structure of the fake-barn case rather than Gettier's: the luck lies in the environment, which makes error easy, not in an intervention between the source and the fact. Pritchard grants that environmental luck is compatible with understanding \citep{Pritchard2010}. The argument extends to intervening luck. Suppose that the book she reads was written by an author who invented his history without research and, by coincidence, got it right. Her belief is then Gettierized in the classical way. But nothing about her grasp has changed: she can answer the same questions, draw the same inferences, and apply the explanation to the same comparisons with other plains peoples. If understanding required safety, she would understand in the first case and fail to understand in the second, though her cognitive relation to the explanation is the same in both. What changes between the cases is the modal profile of her source, and that is what knowledge tracks and understanding does not.

\begin{theorem}\label{t:ideal}
A subject can understand why $p$ through an explanation some of whose peripheral propositions are false, and therefore cannot be known.
\end{theorem}

\begin{proof}
By Definition~\ref{d:trueenough}, an explanation may be true enough while containing false peripheral propositions. A chemistry student who explains the rise in pressure of a heated gas in a closed container by the kinetic theory with point molecules and no intermolecular forces accepts an explanation that is true enough: pressure depends on the rate and momentum of molecular collisions with the walls, both rise with temperature, and the idealizations do not alter this dependence. If she can say what would happen if the container were enlarged or the gas cooled, she grasps the explanation, and by Definition~\ref{d:understand} she understands why the pressure rises. The idealizing propositions are false, and since knowledge entails truth (Definition~\ref{d:know}), they are not known.
\end{proof}

A defender of the reductive view may reply that the student knows the core propositions, and that understanding is knowledge of the core. The reply concedes the point that matters. The student's understanding is achieved by reasoning through the idealized model; her ability to answer the what-if questions is exercised on the model, and the model is not something she knows. To say that understanding is knowledge of the core is to identify understanding with a state that lacks the very abilities that make it understanding, a point to which Theorem~\ref{t:noU} returns.

\section{Knowledge without understanding}

\begin{theorem}\label{t:noU}
Knowledge that $q$ explains $p$, and hence knowledge why $p$, does not entail understanding why $p$.
\end{theorem}

\begin{proof}
A homeowner is told by the fire investigator that her house burned down because of a short circuit in the attic wiring. The investigator knows this, and by the thesis of EPI-08 testimony transmits his knowledge: the homeowner's belief is justified, and safe, since in close worlds in which she believes on the testimony of this investigator the testimony is correct. She therefore knows that the fire was caused by the short circuit, and so knows why the house burned down. But she has no idea how a short circuit causes a fire. She cannot say whether the fire would have occurred had the wiring been insulated differently, had the circuit been fused, or had the attic been empty. She does not grasp the explanation (Definition~\ref{d:grasp}), and by Definition~\ref{d:understand} she does not understand why the house burned down.
\end{proof}

The example is Pritchard's \citep{Pritchard2010}, and similar cases are given by \citet{Hills2016} for moral understanding. It shows that the extra ingredient in understanding cannot be supplied by further items of the same kind as those already known. Suppose the investigator adds that short circuits produce heat, that heat ignites dry timber, and that the attic was full of dry timber. The homeowner now knows more explanatory propositions. Whether she understands depends on whether she can put them together, which is an ability and not one more proposition. A subject who knows every proposition of an explanation and cannot use them to answer any what-if question has the materials of understanding without understanding. The point is Ryle's distinction between knowing that and knowing how applied to explanation \citep{Ryle1949}, though nothing here depends on whether knowing how is itself a kind of knowing that.

\begin{corollary}\label{c:neither}
Understanding why $p$ is neither necessary nor sufficient for knowledge why $p$; understanding is not a species of knowledge.
\end{corollary}

\begin{proof}
By Theorem~\ref{t:luck}, a subject can understand why $p$ without knowing the core of the explanation, and so without knowing why $p$; by Theorem~\ref{t:noU}, a subject can know why $p$ without understanding why $p$. If understanding were a species of knowledge, every case of understanding why $p$ would be a case of knowing that some explanation of $p$ holds.
\end{proof}

\section{Objections and replies}

\begin{objection}[Understanding why is knowing why]\label{o:sliwa}
Ordinary usage treats \enquote{she understands why the house burned down} and \enquote{she knows why the house burned down} as nearly interchangeable, and the deficient cases are cases of minimal understanding, which is all that minimal knowledge why provides. Degrees of understanding are degrees of knowledge of explanatory facts \citep{Sliwa2015}.
\end{objection}

\begin{reply}
The near-interchangeability holds in the central cases, in which a subject knows an explanation and grasps it; these are the cases in which understanding and knowledge coincide, and their frequency explains the usage. It does not survive the cases of Theorems~\ref{t:luck} and~\ref{t:noU}. The homeowner, asked whether she understands why her house burned down, would rightly answer that she knows the cause but does not understand how it did it. On the degree view, the further knowledge needed is knowledge of further explanatory propositions, and we have argued that a subject can know all of them without grasp. The degree view also has the wrong shape for idealization: the student's understanding of the gas increases as she masters the idealized model, though the propositions she comes to accept are false and the number of truths she knows does not increase.
\end{reply}

\begin{objection}[Luck undermines understanding too]\label{o:grimm}
A student who reads the one accurate book among fakes has a lucky understanding. Understanding that could so easily have been misunderstanding is defective, and if it is not defective as understanding, safety is not what defeats knowledge in Gettier cases either \citep{Grimm2006}.
\end{objection}

\begin{reply}
The student's state is defective in one respect: she is in no position to tell that her book is reliable, and in a nearby world she would have misunderstood. This is a defect in her epistemic position with respect to the truth of the explanation, which is what the safety condition registers and what knowledge requires. It is not a defect in her understanding, which consists in a correct explanation grasped. A misunderstanding in the nearby world would be a different state, the grasp of a false explanation, and its nearness does not alter what she has in the actual world. Knowledge is the state whose value lies in non-accidental correctness, which is why it is more stable than true belief (EPI-14); understanding is the state whose value lies in seeing connections, and seeing them by luck is still seeing them.
\end{reply}

\begin{objection}[Objectual understanding]\label{o:object}
The argument treats understanding why. But much understanding is objectual, understanding of a subject matter, such as the history of the Comanche or quantum mechanics. Perhaps objectual understanding is a kind of knowledge even if understanding why is not.
\end{objection}

\begin{reply}
Objectual understanding of a subject matter is the grasp of the explanatory relations among its parts, and on any account it is at least as demanding in grasp as understanding why. The arguments transfer: the student understands the history of the Comanche though her beliefs about it are unsafe, and a reader of an encyclopedia can know many facts about quantum mechanics while understanding none of it. Objectual understanding also tolerates more peripheral falsehood than understanding why, since a subject matter has more periphery, and \citet{Kvanvig2003} argues on this basis that it is less closely tied to truth than knowledge is.
\end{reply}

\begin{objection}[Factivity]\label{o:fact}
If understanding tolerates falsehoods, it is no longer an epistemic success. A theory that is false in its idealizations would then confer understanding on those who grasp it, and so would a false theory that is \enquote{true enough} by some lax standard. Phlogiston theory would give understanding of combustion.
\end{objection}

\begin{reply}
Definition~\ref{d:trueenough} requires the core to be true and allows falsehood only where it leaves the dependence structure intact. Phlogiston theory gets the dependence wrong: it treats combustion as the release of a substance rather than combination with one, and it misrepresents which factors combustion depends on and in which direction mass changes. It therefore fails the definition and gives no understanding of combustion, though it may give understanding of some of its correlations. Understanding is quasi-factive: factive in the core, tolerant in the periphery. That is the right position between Kvanvig's and Elgin's more permissive views and the reductive view, on which every false proposition accepted would be a defect.
\end{reply}

\section{What the result does not establish}

The argument does not show that understanding is more valuable than knowledge, or the reverse. It shows that each has a value the other lacks: knowledge the value of non-accidental truth, understanding the value of grasp. It does not show that understanding requires no justification at all; we have omitted a justification condition because the cases in which understanding seems to require it are cases in which the subject's grasp is defective, and we have not argued that there are none. Nor does the argument address understanding of language, of persons or of works of art, which may have different structures. Finally, Corollary~\ref{c:neither} concerns understanding why and objectual understanding as states of individuals; it is compatible with the claim that knowledge is the normal route to understanding, and that most understanding is accompanied by knowledge of its core.

\section{Conclusion}

Understanding why something is so consists in grasping an explanation of it that is true enough. It requires that the core of the explanation be true but not that the subject's belief in it be safe, and it tolerates falsehood in idealizations and approximations. Knowledge requires safety and truth throughout, but not grasp. Hence a subject can understand without knowing, as when her source is reliable by luck or her model is idealized, and can know without understanding, as when she has the explanation on good testimony without being able to use it. Understanding is not a kind of knowledge.

\begin{thebibliography}{99}\small
\bibitem[Burnyeat(1981)]{Burnyeat1981} Burnyeat, M.~F. (1981). Aristotle on understanding knowledge. In E. Berti (Ed.), \emph{Aristotle on Science: The Posterior Analytics} (pp.~97--139). Antenore.
\bibitem[Elgin(2007)]{Elgin2007} Elgin, C.~Z. (2007). Understanding and the facts. \emph{Philosophical Studies}, 132(1), 33--42.
\bibitem[Elgin(2017)]{Elgin2017} Elgin, C.~Z. (2017). \emph{True Enough}. MIT Press.
\bibitem[Grimm(2006)]{Grimm2006} Grimm, S.~R. (2006). Is understanding a species of knowledge? \emph{The British Journal for the Philosophy of Science}, 57(3), 515--535.
\bibitem[Hills(2016)]{Hills2016} Hills, A. (2016). Understanding why. \emph{No\^us}, 50(4), 661--688.
\bibitem[Khalifa(2017)]{Khalifa2017} Khalifa, K. (2017). \emph{Understanding, Explanation, and Scientific Knowledge}. Cambridge University Press.
\bibitem[Kvanvig(2003)]{Kvanvig2003} Kvanvig, J.~L. (2003). \emph{The Value of Knowledge and the Pursuit of Understanding}. Cambridge University Press.
\bibitem[Pritchard(2010)]{Pritchard2010} Pritchard, D. (2010). Knowledge and understanding. In D. Pritchard, A. Millar and A. Haddock, \emph{The Nature and Value of Knowledge: Three Investigations}. Oxford University Press.
\bibitem[Ryle(1949)]{Ryle1949} Ryle, G. (1949). \emph{The Concept of Mind}. Hutchinson.
\bibitem[Sliwa(2015)]{Sliwa2015} Sliwa, P. (2015). Understanding and knowing. \emph{Proceedings of the Aristotelian Society}, 115(1), 57--74.
\bibitem[Woodward(2003)]{Woodward2003} Woodward, J. (2003). \emph{Making Things Happen: A Theory of Causal Explanation}. Oxford University Press.
\bibitem[Zagzebski(2001)]{Zagzebski2001} Zagzebski, L. (2001). Recovering understanding. In M. Steup (Ed.), \emph{Knowledge, Truth, and Duty: Essays on Epistemic Justification, Responsibility, and Virtue} (pp.~235--251). Oxford University Press.
\end{thebibliography}

\end{document}
